% GENERATED FILE. Edit canonical lessons, then run npm run generate:books.
% canonical-source-hash: fnv1a64:7e9d14d98b210f1c
% canonical-lessons: RU-C341-protivorechit, RU-C341-otvergat, RU-C341-kompromiss, RU-C341-nastaivat, RU-C341-priznavat

\chapter{Reaching a compromise}
\label{ch:ru-reaching-a-compromise}
\hlchaptermodality{341}

\begin{chapteropening}
\textbf{By the end of this chapter:} \emph{I can say that two things contradict each other, reject an idea, insist on something, admit a point and reach a compromise.}

Complete the last lesson of chapter 341: I can say that two things contradict each other, reject an idea, insist on something, admit a point and reach a compromise.
\end{chapteropening}

\section[\texorpdfstring{\ru{противоречить} (protivoréchit)}{противоречить (protivoréchit)}]{\ru{противоречить} (protivoréchit) — to contradict}

\label{lesson:RU-C341-protivorechit}



\begin{quote}
Before the new one: say the Russian for to object, then the Russian for an objection.
\end{quote}

\subsection*{You'll want to know: \ru{противоречить}}
\textbf{\ru{противоречить}} (\emph{protivoréchit}) — \textquotedblleft{}to contradict\textquotedblright{}.

\begin{quote}
\ru{Ты} \ru{сам} \ru{себе} \ru{противоречишь}.
\end{quote}

\begin{quote}
\emph{Ty sam sebé protivoréchish.}
\end{quote}

\begin{quote}
You are contradicting yourself.
\end{quote}

Imperfective only. It takes the dative: \emph{\ru{противоречить} \ru{фактам}}.

\begin{cousinweb}
\textbf{\ru{противоречить}} is built from \emph{\ru{противо}-} \textquotedblleft{}against\textquotedblright{} and \emph{\ru{речь}} \textquotedblleft{}speech\textquotedblright{}, which is inherited from Proto-Slavic \emph{*rěčĭ}, from \emph{*rekti} \textquotedblleft{}to say\textquotedblright{} (EtymDB). To contradict is to speak against.
\end{cousinweb}

\begin{grammarlens}[title={What you've built}]
One more word for telling, explaining and describing at B1.
\end{grammarlens}

\subsection*{Your turn}
\begin{itemize}
\raggedright
  \item \emph{Say it:} \emph{protivoréchit}
  \item \emph{Say it:} \emph{protivoréchit}, once more
  \item \emph{Say it:} a sentence of your own with \emph{protivoréchit}
  \item \emph{Recall:} say the Russian for to object, then the Russian for an objection, then say \emph{protivoréchit} again
\end{itemize}

\begin{tcolorbox}[breakable,colback=teal!4,colframe=teal!35!black,title={Before you move on}]
What does \ru{противоречить} mean? (\textquotedblleft{}To contradict\textquotedblright{}.) Say it once more.
\end{tcolorbox}
\section[\texorpdfstring{\ru{отвергать} (otvergát)}{отвергать (otvergát)}]{\ru{отвергать} (otvergát) — to reject}

\label{lesson:RU-C341-otvergat}



\begin{quote}
Before the new one: say the Russian for an objection, then the Russian for to contradict.
\end{quote}

\subsection*{You'll want to know: \ru{отвергать}}
\textbf{\ru{отвергать}} (\emph{otvergát}) — \textquotedblleft{}to reject\textquotedblright{}.

\begin{quote}
\ru{Начальник} \ru{отверг} \ru{наше} \ru{предложение}.
\end{quote}

\begin{quote}
\emph{Nachálnik otvérg náshe predlozhéniye.}
\end{quote}

\begin{quote}
The boss rejected our proposal.
\end{quote}

Imperfective; perfective \emph{\ru{отвергнуть}} (past \emph{\ru{отверг}, \ru{отвергла}}). Formal; in conversation people say \emph{\ru{не} \ru{принять}}.

\begin{cousinweb}
\textbf{\ru{отвергать}} is built from \emph{\ru{от}-} \textquotedblleft{}away\textquotedblright{} and \emph{\ru{вергать}} \textquotedblleft{}to throw\textquotedblright{}, a bookish verb. EtymDB traces \emph{\ru{отвергнуть}} to Proto-Slavic \emph{*vĭrgnǫti} (EtymDB), usually glossed \textquotedblleft{}to throw\textquotedblright{}. To reject is to throw away.
\end{cousinweb}

\begin{grammarlens}[title={What you've built}]
One more word for telling, explaining and describing at B1.
\end{grammarlens}

\subsection*{Your turn}
\begin{itemize}
\raggedright
  \item \emph{Say it:} \emph{otvergát}
  \item \emph{Say it:} \emph{otvergát}, once more
  \item \emph{Say it:} a sentence of your own with \emph{otvergát}
  \item \emph{Recall:} say the Russian for an objection, then the Russian for to contradict, then say \emph{otvergát} again
\end{itemize}

\begin{tcolorbox}[breakable,colback=teal!4,colframe=teal!35!black,title={Before you move on}]
What does \ru{отвергать} mean? (\textquotedblleft{}To reject\textquotedblright{}.) Say it once more.
\end{tcolorbox}
\section[\texorpdfstring{\ru{компромисс} (kompromíss)}{компромисс (kompromíss)}]{\ru{компромисс} (kompromíss) — a compromise}

\label{lesson:RU-C341-kompromiss}



\begin{quote}
Before the new one: say the Russian for to contradict, then the Russian for to reject.
\end{quote}

\subsection*{You'll want to know: \ru{компромисс}}
\textbf{\ru{компромисс}} (\emph{kompromíss}) — \textquotedblleft{}a compromise\textquotedblright{}.

\begin{quote}
\ru{После} \ru{долгого} \ru{спора} \ru{мы} \ru{нашли} \ru{компромисс}.
\end{quote}

\begin{quote}
\emph{Pósle dólgogo spóra my nashlí kompromíss.}
\end{quote}

\begin{quote}
After a long argument we found a compromise.
\end{quote}

Masculine. \emph{\ru{Пойти} \ru{на} \ru{компромисс}} — to make a compromise.

\begin{cousinweb}
\textbf{\ru{компромисс}} is usually traced to French \emph{compromis}, from Latin \emph{compromissum} \textquotedblleft{}a mutual promise to accept what an arbiter decides\textquotedblright{} (Lewis \& Short), built on \emph{mittere} \textquotedblleft{}to send\textquotedblright{}.
\end{cousinweb}

\begin{grammarlens}[title={What you've built}]
One more word for telling, explaining and describing at B1.
\end{grammarlens}

\subsection*{Your turn}
\begin{itemize}
\raggedright
  \item \emph{Say it:} \emph{kompromíss}
  \item \emph{Say it:} \emph{kompromíss}, once more
  \item \emph{Say it:} a sentence of your own with \emph{kompromíss}
  \item \emph{Recall:} say the Russian for to contradict, then the Russian for to reject, then say \emph{kompromíss} again
\end{itemize}

\begin{tcolorbox}[breakable,colback=teal!4,colframe=teal!35!black,title={Before you move on}]
What does \ru{компромисс} mean? (\textquotedblleft{}A compromise\textquotedblright{}.) Say it once more.
\end{tcolorbox}
\section[\texorpdfstring{\ru{настаивать} (nastáivat)}{настаивать (nastáivat)}]{\ru{настаивать} (nastáivat) — to insist}

\label{lesson:RU-C341-nastaivat}



\begin{quote}
Before the new one: say the Russian for to reject, then the Russian for a compromise.
\end{quote}

\subsection*{You'll want to know: \ru{настаивать}}
\textbf{\ru{настаивать}} (\emph{nastáivat}) — \textquotedblleft{}to insist\textquotedblright{}.

\begin{quote}
\ru{Врач} \ru{настаивает}, \ru{чтобы} \ru{я} \ru{больше} \ru{спал}.
\end{quote}

\begin{quote}
\emph{Vrach nastáivayet, chtóby ya bólshe spal.}
\end{quote}

\begin{quote}
The doctor insists that I sleep more.
\end{quote}

Imperfective; perfective \emph{\ru{настоять}}. \emph{\ru{Настаивать} \ru{на}} + prepositional, or \emph{\ru{настаивать}, \ru{чтобы}} + past tense.

\begin{cousinweb}
\textbf{\ru{настаивать}} is built from \emph{\ru{на}-} \textquotedblleft{}on\textquotedblright{} and \emph{\ru{стоять}} \textquotedblleft{}to stand\textquotedblright{}, inherited from Proto-Slavic \emph{*stojati} (EtymDB). To insist is to stand on your point.
\end{cousinweb}

\begin{grammarlens}[title={What you've built}]
One more word for telling, explaining and describing at B1.
\end{grammarlens}

\subsection*{Your turn}
\begin{itemize}
\raggedright
  \item \emph{Say it:} \emph{nastáivat}
  \item \emph{Say it:} \emph{nastáivat}, once more
  \item \emph{Say it:} a sentence of your own with \emph{nastáivat}
  \item \emph{Recall:} say the Russian for to reject, then the Russian for a compromise, then say \emph{nastáivat} again
\end{itemize}

\begin{tcolorbox}[breakable,colback=teal!4,colframe=teal!35!black,title={Before you move on}]
What does \ru{настаивать} mean? (\textquotedblleft{}To insist\textquotedblright{}.) Say it once more.
\end{tcolorbox}
\section[\texorpdfstring{\ru{признавать} (priznavát)}{признавать (priznavát)}]{\ru{признавать} (priznavát) — to admit, to acknowledge}

\label{lesson:RU-C341-priznavat}



\begin{quote}
Before the new one: say the Russian for a compromise, then the Russian for to insist.
\end{quote}

\subsection*{You'll want to know: \ru{признавать}}
\textbf{\ru{признавать}} (\emph{priznavát}) — \textquotedblleft{}to admit, to acknowledge\textquotedblright{}.

\begin{quote}
\ru{Он} \ru{не} \ru{любит} \ru{признавать} \ru{свои} \ru{ошибки}.
\end{quote}

\begin{quote}
\emph{On ne lyúbit priznavát svoí oshíbki.}
\end{quote}

\begin{quote}
He does not like admitting his mistakes.
\end{quote}

Imperfective; present \emph{\ru{признаю}, \ru{признаёшь}}; perfective \emph{\ru{признать}}. Built on \emph{\ru{знать}}.

\begin{cousinweb}
\textbf{\ru{признавать}} is built from \emph{\ru{при}-} and \emph{\ru{знать}} \textquotedblleft{}to know\textquotedblright{}, which is inherited from Proto-Slavic \emph{*znati} (EtymDB). To admit a thing is to own that you know it.
\end{cousinweb}

\begin{grammarlens}[title={What you've built}]
One more word for telling, explaining and describing at B1.
\end{grammarlens}

\subsection*{Your turn}
\begin{itemize}
\raggedright
  \item \emph{Say it:} \emph{priznavát}
  \item \emph{Say it:} \emph{priznavát}, once more
  \item \emph{Say it:} a sentence of your own with \emph{priznavát}
  \item \emph{Recall:} say the Russian for a compromise, then the Russian for to insist, then say \emph{priznavát} again
\end{itemize}

\begin{tcolorbox}[breakable,colback=teal!4,colframe=teal!35!black,title={Before you move on}]
What does \ru{признавать} mean? (\textquotedblleft{}To admit, to acknowledge\textquotedblright{}.) Say it once more.
\end{tcolorbox}
